\subsection{An explicit finite target retaining the entire central defect}

Put $m=2^{14}3^6=|E|$, and retain the actual open normal
free subgroup $L_0=\ker(B\to\mathcal T)$ of rank $25$.
Define the actual finite target and its induced source by
\[
\begin{aligned}
 V_L&=\overline{\gamma_3(L_0)L_0^{,2m}},&
 Q_*&=B/V_L,\\
 R_*&=U/(V_L)^6.&
\end{aligned}                                                   \tag{B.S12}
\]
Here $L_0^{,2m}$ denotes the closed subgroup generated
by all $2m$-th powers.  It is characteristic in $L_0$,
so $V_L$ is normal in $B$.  Every positive target action
preserving the cusp class preserves it.  On the pure
subgroup of $U$, the six literal maps give exactly
\[
 R_*=B/\operatorname{core}_B
             \ker(q_*\varphi_P:F\longrightarrow Q_*),           \tag{B.S13}
\]
where $q_*:B\twoheadrightarrow Q_*$ is the quotient in
(B.S12).  This source kernel is contained in $B'$,
since its image in $U$ is contained in $L_0^6\subseteq J'$.
Thus no additional abelian source factor is needed.

\begin{result}{B.2.17}{Proved}{A finite canonical cover detects the full central kernel}
The natural quotient $U^{\mathrm c}\twoheadrightarrow
R_*^{\mathrm c}$ is injective on its actual central
kernel $M=\ker(U^{\mathrm c}\to U)$.  The displayed finite
groups have exact orders
\[
 |Q_*|=2^{4579}3^{1952},\qquad
 |R_*|=2^{27464}3^{11706}.                                      \tag{B.S14}
\]
There is also an explicitly defined characteristic finite
target refining $Q_*$, whose induced canonical source
cover has the same injectivity property.
\end{result}

\begin{proof}
Let $N=L_0^6\triangleleft U$, and let $P$ be its full
inverse image in $U^{\mathrm c}$.  Its kernel over $N$
is the central finite group $M$, annihilated by $m$.
Each factor $L_0$ of $N$ is free profinite.  Lifting its
free basis to $P$ gives a continuous section of that
factor, since its composite back to $L_0$ is the identity
on the free basis.  Choose one such section for each of
the six factors.  Every element of $P$ then has a unique
expression
\[
                 z\,s_1(v_1)\cdots s_6(v_6),
                 \qquad z\in M,\quad v_i\in L_0,
\]
with the factor order fixed.  The commutator of two
different section factors lies in $M$, hence is central.
It defines a continuous bilinear map from their two
abelianizations to $M$, and factors through reduction
of those abelianizations modulo $m$.

For each factor take its free class-two quotient of
exponent $2m$,
$C=L_0/\overline{\gamma_3(L_0)L_0^{,2m}}$.
It is finite.  An explicit collected form gives
$(\mathbb Z/2m)^{25}$ for the generator exponents and
$(\mathbb Z/m)^{300}$ for the commutator exponents.
Indeed the class-two power formula forces commutator
exponent $m$ when all elements have exponent $2m$;
conversely these collected forms with the standard
bilinear multiplication satisfy the relations and
give every element uniquely.  Thus
\[
                         |C|=(2m)^{25}m^{300}.                 \tag{B.S15}
\]

The expression for $P$ gives a finite group on the set
$M\times C^6$: retain the multiplications in the six
copies of $C$, and use the original cross-factor
bilinear commutators valued in $M$ when collecting
the factors in their fixed order.  These commutators
are defined on the quotient because they depend only
on the generator exponents modulo $m$.  Associativity
is the bilinearity identity.  The group has class at
most two and exponent dividing $2m$: the extra central
term in a $2m$-th power is a multiple of
$\binom{2m}{2}=m(2m-1)$ and hence vanishes in $M$.
This supplies an actual homomorphism from $P$ to that
finite group, which is the identity on its subgroup $M$.
Consequently, with
\[
               V_P=\overline{\gamma_3(P)P^{,2m}},
               \qquad M\cap V_P=1.                            \tag{B.S16}
\]
The subgroup $V_P$ is characteristic in $P$ and thus
normal in $U^{\mathrm c}$.  Its image in $N$ is
$\gamma_3(N)N^{,2m}=V_L^6$.  Therefore the inverse
image in $U^{\mathrm c}$ of the kernel of $U\to R_*$
is exactly $MV_P$.  The canonical quotient kernel is
\[
 \ker(U^{\mathrm c}\longrightarrow R_*^{\mathrm c})
       =\overline{[U^{\mathrm c},MV_P]}
       =\overline{[U^{\mathrm c},V_P]}\subseteq V_P.             \tag{B.S17}
\]
The first equality is precisely the actual commutator
comparison (B.S3).  Equations (B.S16)--(B.S17) prove
the asserted injectivity on $M$.

The target and source orders follow from (B.S15),
$[B:L_0]=144$, and $[U:L_0^6]=m$:
\[
 |Q_*|=144(2m)^{25}m^{300},\qquad
 |R_*|=m\bigl((2m)^{25}m^{300}\bigr)^6.
\]
Substitution of $m=2^{14}3^6$ gives (B.S14).

Finally let $n=|Q_*|$, and define the open characteristic
subgroup
\[
 K^{\#}=\bigcap_{\substack{W\triangleleft_{\rm o}B\\
                           [B:W]\leq n}}W,
             \qquad Q^{\#}=B/K^{\#}.                           \tag{B.S18}
\]
The intersection is finite because $B$ is finitely
generated and has only finitely many open subgroups
of each bounded index.  It includes $V_L$ among its
intersected subgroups, so $Q^{\#}$ maps to $Q_*$ by
the actual quotient map.  Take the induced source
$E_{Q^{\#}}$ as in (B.Q5).  Its canonical cover maps
to $R_*^{\mathrm c}$.  The composite of the natural
map from $U^{\mathrm c}$ with this map is the one in
(B.S17), proving injectivity on $M$ at this
characteristic target as well.
\end{proof}

This result supplies an explicit finite detector for the
whole central kernel.  Its orders are bounds of a usable
mathematical construction, not claims that the resulting
groups have been enumerated in this audit.  The complete
positive framed-automorphism census at this target has
not been performed.  Thus the actual finite-index result
B.2.16 is retained, while the index-one question remains
unresolved.
